\section{引子}

以掷一颗骰子为例，若记事件：
\begin{align*}
&A_1=\left\{ \text{点数比}1\text{小} \right\} \\
&A_2=\left\{ \text{点数比}2\text{小} \right\} \\
&A_3=\left\{ \text{点数比}3\text{小} \right\} \\
&A_4=\left\{ \text{点数比}4\text{小} \right\} \\
&A_5=\left\{ \text{点数比}5\text{小} \right\} \\
&A_6=\left\{ \text{点数比}6\text{小} \right\}
\end{align*}
它们的概率可用图形显示如下左图。
很明显，事件概率依次递增，且呈现等差数列的规律。

\begin{figure}[h]
\centering
\includegraphics[height=2.5cm]{2.1.png}
\end{figure}

若记事件：
\begin{align*}
&A_1=\left\{ \text{点数能被}1\text{整除} \right\} \\
&A_2=\left\{ \text{点数能被}2\text{整除} \right\} \\
&A_3=\left\{ \text{点数能被}3\text{整除} \right\} \\
&A_4=\left\{ \text{点数能被}4\text{整除} \right\} \\
&A_5=\left\{ \text{点数能被}5\text{整除} \right\} \\
&A_6=\left\{ \text{点数能被}6\text{整除} \right\}
\end{align*}
则它们的概率用图形显示如上右图。
事件概率呈现依次递减的规律。

通过图形，我们可以直观地感受事件概率的分布规律，这就是本章要讨论的内容。




